\documentclass[10pt,letterpaper]{article}

\usepackage[letterpaper,left=0.65in,right=0.65in,top=0.35in,bottom=0.4in]{geometry}
\usepackage{fontspec}
\usepackage{enumitem}
\usepackage{xcolor}
\definecolor{linkblue}{HTML}{234A70}
\usepackage[colorlinks=true,urlcolor=linkblue]{hyperref}
\usepackage[normalem]{ulem}

\setmainfont{Arial}
\pagestyle{empty}
\setlength{\parindent}{0pt}
\setlength{\parskip}{0pt}
\raggedright
\urlstyle{same}
\hypersetup{
  pdftitle={Javier Contreras Tenorio - Resume},
  pdfauthor={Javier Contreras Tenorio}
}

\setlist[itemize]{
  leftmargin=14pt,
  labelsep=5pt,
  topsep=1pt,
  partopsep=0pt,
  itemsep=1pt,
  parsep=0pt
}
\setlist[itemize,1]{label=\textbullet}
\setlist[itemize,2]{label=\textendash}

\newcommand{\resumesection}[1]{%
  \par\addvspace{5pt}%
  {\fontsize{12}{13}\selectfont #1\par}%
  \vspace{0.5pt}\hrule height 0.7pt\vspace{3pt}%
}
\newcommand{\role}[2]{%
  \par\addvspace{6pt}%
  \textbf{#1}\hfill\textit{#2}\par\nobreak
}
\newcommand{\resumelink}[2]{\href{#1}{\uline{#2}}}

\begin{document}
\fontsize{10}{12}\selectfont

{\centering
  {\fontsize{12}{14}\selectfont\textbf{Javier Contreras Tenorio}\par}
  \vspace{3pt}
  (734) 596 0204 - \resumelink{mailto:javierct099@gmail.com}{javierct099@gmail.com}
  – LinkedIn: \resumelink{http://www.linkedin.com/in/javier-contreras-t}{www.linkedin.com/in/javier-contreras-t}
  – \resumelink{https://jav099.github.io/}{jav099.github.io}\par
}

\resumesection{Education}
\makebox[0.6\linewidth][l]{\textbf{UNIVERSITY OF MICHIGAN}, Ann Arbor, United States.}2018 - 2022\par
\begin{itemize}
  \item B.S. Computer Science – 3.623 GPA
\end{itemize}
\makebox[0.6\linewidth][l]{\textbf{UNIVERSITY OF WASHINGTON}, Seattle, United States}2025 - Present\par
\begin{itemize}
  \item M.S. Computer Science – 4.0 GPA
\end{itemize}

\resumesection{Skills}
\textit{\textbf{C++,} Python, Mentoring, Algorithms, Browser Architecture, Performance, Code Reviewing, Agentic tools}\par

\resumesection{Experience}
\role{SOFTWARE ENGINEER I\\SOFTWARE ENGINEER II, MICROSOFT, SEATTLE}{August 2022 - Present}
Part of Microsoft Edge’s Accessibility and Layout team for the \textbf{Chromium} web platform
\begin{itemize}
  \item Designed and implemented fixes for bugs with Chromium and Edge’s Accessibility code
  \begin{itemize}
    \item Helped Microsoft Edge stay accessibility compliant by making sure that any accessibility related bugs were fixed within the allowed window.
    \item This involved debugging multiple processes and implementing multi-process solutions.
  \end{itemize}
  \item Designed and implemented part of the ViewsAX project for \textbf{Chromium}.
  \begin{itemize}
    \item This project consisted in refactoring Chromiums’ Views’ \textbf{accessibility} architecture to more conform to Blink’s existing architecture so the existing mechanisms in the browser process could be used to serialize and unserialize each View’s accessibility information.
    \item Took on a lead role in certain sub-projects of the overarching project by making design documents for certain systems for the project, as well as onboarding new members working on the project and reviewing their code.
  \end{itemize}
  \item Became \textbf{owner} of //ui/views/accessibility directory in Chromium.
  \item Contributed over 200 \textbf{code reviews} and over 200 changes to Chromium (under the account \resumelink{mailto:javiercon@microsoft.com}{\textbf{javiercon@microsoft.com}})
  \item Shipped \textbf{CSS Gap Decorations} project, new CSS properties in Chromium browsers.
  \begin{itemize}
    \item Implemented the feature end to end, from the Layout side of the browser all the way to the painting side.
    \item Helped \textbf{write and publish the official spec} for this CSS feature, and wrote hundreds of WPTs for the feature
    \item Collaborated with Google, the CSS Working Group, and web developers to shape the feature.
  \end{itemize}
  \item Developed internal tools, skills, agents to improve Edge’s and Edge Web Platform’s agentic workflows.
  \begin{itemize}
    \item Developed an internal autonomous dev workflow engine that turns Copilot CLI into a task queue and pipeline orchestrator for reusable pipelines of AI prompts, scripts etc, for Copilot to execute autonomously.
  \end{itemize}
\end{itemize}

\role{TOP MICROSOFT CHROMIUM CONTRIBUTOR}{August 2022 - Present}
Since starting at Microsoft in August 2022, became the top MSFT contributor to Chromium (commit count) during that period.\par
My work spanned across accessibility, CSS, layout, animations, painting, and Chromium's agentic tooling (skills, instructions, etc.)

\role{INTERN MENTOR, MICROSOFT}{Summer 2026}
Mentored an intern and helped them ship a CSS feature in Chromium.

\role{SOFTWARE ENGINEER INTERNSHIP, MICROSOFT (Remote)}{May 2021 - July 2021}

\resumesection{Outreach \& Volunteering}
\role{SPEAKER, BLINKON CONFERENCE}{April 2026}
Presented CSS Gap Decorations at Chromium’s BlinkOn conference

\role{INVITED SPEAKER, UNIVERSITY OF MICHIGAN}{October 2026}
Gave a talk about contributing to open source and Chromium at the University of Michigan.

\role{TEALS VOLUNTEER, SEATTLE PUBLIC SCHOOLS}{Sept 2023 - May 2025}
Volunteered as a Computer Science high school teacher at Seattle World School
\begin{itemize}
  \item Designed and planned out introductory computer science lessons in English and Spanish
  \item Taught these lessons to high school kids, most of whom are not very proficient in English since it is a school for immigrants.
\end{itemize}

\end{document}
